\section{范畴迹的构造}
近年来，范畴中心/范畴迹的概念在表示论中扮演了重要角色。
关于它们在特征层理论中的应用，我们参看 \cite{BN,BFO,Lu3,Lu4,Lu5,Lu6}；关于在稳定 Bernstein 中心中的应用，参看 \cite{BKV}。本文余下部分基于与肖梁的合作工作 \cite{XZ}，说明这些范畴构造的另一种视角与应用。

\subsection{范畴迹}
\label{S: cat trace}
为说明定义的动机，我们先简要回顾代数的迹（有时也称为余中心）的概念。
设 $E$ 是基交换环，$A$ 是 $E$-代数。从 $A$ 到 $E$-模 $V$ 的迹函数指的是满足 $f(ab)=f(ba)$ 的 $E$-线性映射 $f: A\to V$。
$A$ 的迹是一个 $E$-模 $\Tr(A)$，配以迹函数 $\tr: A\to \Tr(A)$，使得每个迹函数 $f:A\to V$ 都唯一地分解为 $A\xrightarrow{\tr} \Tr(A)\xrightarrow{\bar{f}}V$，其中 $\bar{f}:\Tr(A)\to V$ 是某个线性映射。显然，若把 $\Tr(A)$ 定义为 $A$ 模去由形如 $ab-ba$（$a,b\in A$）的元素张成的 $E$-子模所得的商，并令 $\tr$ 为该商映射，则 $(\Tr(A),\tr)$ 就是 $A$ 的迹。等价地，也可以把 $\Tr(A)$ 表示成 $E$-模范畴中下图之余极限
\begin{equation}
\label{E:HH}
\Tr(A)=\on{colim}(A\otimes A\substack{\longrightarrow\\[-1em] \longrightarrow}  A)
\end{equation}
其中两个映射 $A\otimes A\to A$ 分别由 $a\otimes b\mapsto ab$ 与 $a\otimes b\mapsto ba$ 给出。注意这个图正是计算 $A$ 的 Hochschild 同调的 bar 复形的前两项。

下面转到范畴情形，引入幺半范畴的范畴迹概念。
基本想法如下：幺半范畴是某类范畴构成的对称幺半 $2$-范畴中的代数对象；一般地，应当把对称幺半（高阶）范畴中代数对象的迹定义为其 bar 复形 \eqref{E:HH} 的余极限。借助稳健的高阶范畴论，这一想法在相当一般的设定下都可行。例如，我们参看 \cite{BN} 中精彩的论述。

对我们后来的应用而言，只需在通常的 $E$-线性范畴层面引入这一概念，并且重要的是显式地构造范畴迹。另一方面，我们还需要幺半自函子的扭范畴迹的概念。因此我们详细展开定义。我们将在本质小的 $E$-线性范畴构成的 $2$-范畴中工作。
设 $(\mC,\otimes,\mathbf{1})$ 是 $E$-线性幺半范畴，$\sigma: \mC\to \mC$ 是幺半自函子。对每个 $n$，记 $\mC^{\otimes n}$ 为这样的 $E$-线性范畴：对象是 $(X_i)_{i=1,\ldots,n}$，态射为 $\Hom_{\mC^{\otimes n}}((X_i)_i, (Y_i)_i)=\otimes_i \Hom_\mC(X_i,Y_i)$。则左 $\sigma$-扭范畴迹 $\Tr_\sigma(\mC)$ 定义为下图之 $2$-余极限
\[
\mC^{\otimes 3}\ \substack{\longrightarrow\\[-1em] \longrightarrow\\[-1em] \longrightarrow} \ \mC^{\otimes 2}\ \substack{\longrightarrow\\[-1em] \longrightarrow} \ \mC,
\]
其中函子 $\mC^{\otimes 2}\to \mC$ 分别由 $(X,Y)\mapsto X\otimes Y$ 与 $Y\otimes \sigma X$ 给出，而 $\mC^{\otimes 3}\to \mC^{\otimes 2}$ 分别由 $(X,Y,Z)\mapsto (X,Y\otimes Z)$、$(Y, Z\otimes \sigma X)$ 与 $(X\otimes Y,Z)$ 给出。

更具体地，设 $\mD$ 是（通常的）$E$-线性范畴。
从 $\mC$ 到 $\mD$ 的左 $\sigma$-扭\emph{迹函子}指的是一个函子 $F:\mC\to\mD$，连同典范同构族（对每个变元都是函子性的）
\[\al_{X,Y}: F(X\otimes Y)\cong F(Y\otimes \sigma X)\]
使得下图交换
\begin{equation}
\label{E:cond trace}
\xymatrix{
F((X\otimes Y)\otimes Z)\ar^{\al_{X\otimes Y,Z}}[r]\ar_\cong[d] & F(Z\otimes \sigma (X\otimes Y))\ar^{\cong}[r]&  F((Z\otimes \sigma X)\otimes \sigma Y)  \\
F((X\otimes (Y\otimes Z))\ar^{\al_{X,Y\otimes Z}}[r] & F((Y\otimes Z)\otimes \sigma X)\ar^{\cong}[r]&  F(Y\otimes (Z\otimes \sigma X)). \ar_{\al_{Y,(Z\otimes\sigma X)}}[u]
}
\end{equation}
注意由此可推出下面两个图也交换
\begin{equation}
\label{E:cond trace cor}
\xymatrix{
F(\mathbf{1}\otimes X)\ar_{\cong}[d]\ar^{\al_{\mathbf{1},X}}[r]& F(X\otimes \sigma\mathbf{1})\ar^\cong[d]&  F(X\otimes Y)\ar^{\al_{Y,\sigma X}\al_{X,Y}}[r]\ar_\cong[d]& F(\sigma X\otimes \sigma Y)\ar^\cong[d]\\
 F(X)\ar^-\cong[r]& F(X\otimes \mathbf{1})&  F((X\otimes Y)\otimes\mathbf{1})\ar^{\al_{X\otimes Y,\mathbf{1}}} [r]& F(\mathbf{1}\otimes\sigma(X\otimes Y)).
}\end{equation}
记 $\on{Func}^{\tr_\sigma}(\mC,\mD)$ 为左 $\sigma$-扭迹函子构成的范畴。则 $\Tr_\sigma(\mC)$（若存在）是满足下述性质的唯一范畴（在同构意义下唯一，且同构唯一）：对任意（通常的）$E$-线性范畴 $\mD$，
\[\on{Func}^{\tr_\sigma}(\mC,\mD)\cong \on{Func}(\Tr_\sigma(\mC),\mD),\]
其中 $\on{Func}(-,-)$ 表示 $E$-线性范畴之间的 $E$-线性函子范畴。

对我们的应用而言，需要 $\Tr_\sigma(\mC)$ 的同调空间的显式描述。通常，描述范畴余极限图的同调空间很困难。但在我们的情形，对 $\mC$ 作一点温和的假设即可做到。
\begin{prop}
\label{L:trace cat}
设 $\mC$ 是小 $E$-线性范畴，且其中每个对象都有左对偶对象。则 $(\Tr_\sigma(\mC),\tr)$ 存在。若 $\mC$ 是加性的，则 $\Tr_\sigma(\mC)$ 亦然。
\end{prop}
回忆在幺半范畴中，对象 $V$ 的左对偶指的是对象 $V^*$ 连同态射
$$\coev_V:\mathbf{1}\to V\otimes V^*,\quad \ev_V: V^*\otimes V\to \mathbf{1}$$
使得下面两个复合都等于恒等映射
\[V\cong \mathbf{1}\otimes V\xrightarrow{\coev_V\otimes\id_{V}}(V\otimes V^*)\otimes V\cong V\otimes (V^*\otimes V)\xrightarrow{\id_{V}\otimes\ev_V}V\otimes\mathbf{1}\cong V\]
\[V^*\cong V^*\otimes\mathbf{1}\xrightarrow{\id_{V^*}\otimes\coev_V} V^*\otimes (V\otimes V^*)\cong (V^*\otimes V)\otimes V^*\xrightarrow{\ev_V\otimes\id_{V^*}}\mathbf{1}\otimes V^*\cong V^*.\]
三元组 $(V^*,\coev_V,\ev_V)$ 在同构唯一的意义下唯一。
